When an autonomous researcher is evaluated under a fixed wall-clock budget, a faster program also trains on more data before it is scored, so the resource contract can bias what the search appears to discover.
We examine this in AutoPareto, an autonomous Transformer researcher that edits its own training source under three measured objectives: validation quality, cached-decode throughput, and allocated inference memory on a single NVIDIA A40.
After the agent produced its source-edit lineage, we froze seven archived recipes and retrained each on three unseen seeds at exactly matched token exposure.
The discovery-stage quality advantage of the selected agent recipe does not persist under this verification, so it should not be read as an exposure-controlled quality improvement.
Its inference-efficiency advantages do persist: at matched tokens, updates, and batch size, the recipe retains \DecodeRatio$\times$ cached-decode throughput with \ParameterReduction\% fewer parameters and \MemoryReduction\% less allocated inference memory.
Resource contracts and exposure-controlled verification therefore belong in the methodology of autonomous ML research, not only in its reporting.
